% Generated draft: requires mathematical and source review.
\documentclass[11pt]{article}
\usepackage[margin=25mm]{geometry}
\usepackage{fontspec}
\setmainfont{Latin Modern Roman}
\usepackage{amsmath,amssymb,mathtools,unicode-math}
\setmathfont{Latin Modern Math}
\usepackage{enumitem,needspace,xurl,hyperref}
\setlength{\parindent}{0pt}
\setlength{\parskip}{5pt}
\setlength{\emergencystretch}{3em}
\newcommand{\N}{\mathbb{N}}
\newcommand{\Z}{\mathbb{Z}}
\newcommand{\Q}{\mathbb{Q}}
\newcommand{\R}{\mathbb{R}}
\newcommand{\C}{\mathbb{C}}
\newcommand{\F}{\mathbb{F}}
\newcommand{\A}{\mathbb{A}}
\newcommand{\PP}{\mathbb{P}}
\newcommand{\E}{\mathbb{E}}
\newcommand{\eps}{\varepsilon}
\newcommand{\dd}{\,\mathrm{d}}
\begin{document}
\section*{UnsolvedMath ID: 30002479}
\textbf{Source problem code:} catalogue entry
\subsection{Definitions and mathematical statement}
Let \(\kappa\) denote the cardinal parameter in the supplied record. A quasi-order on a set \(X\) is a reflexive, transitive relation \(R\subseteq X\times X\). For quasi-orders \(R\) on \(X\) and \(S\) on \(Y\), a reduction is a map \(f:X\to Y\) satisfying
\[
 x\,R\,x'\quad\Longleftrightarrow\quad f(x)\,S\,f(x')
 \qquad(x,x'\in X).
\]
Borel reducibility additionally requires specified Borel structures on \(X,Y\) and a Borel map \(f\). This definition alone does not define the stronger property called strong invariant universality in the source.

For a specified signature, logic, and sentence \(\varphi\), write \(\operatorname{Mod}_{\kappa}(\varphi)\) for the structures of cardinality \(\kappa\) satisfying \(\varphi\). A descriptive-set-theoretic treatment additionally requires a coding and topology. The supplied record does not specify these data, the sentence \(\varphi\), the definition of generalized tree, the intended embeddings, or the precise meanings of \(\kappa\)-analytic and strongly invariantly universal. Consequently, the following preserves the source questions but remains a conditional formulation requiring those definitions and hypotheses from the cited report.

\begin{enumerate}
\item For the intended class of generalized trees of cardinality \(\kappa\), is the hypothesis printed as
\[
\kappa^{\kappa}=\kappa
\]
necessary in the referenced theorem asserting that embeddability is strongly invariantly universal for \(\kappa\)-analytic quasi-orders?

There is an unresolved defect in this printed hypothesis: for every infinite cardinal \(\kappa\),
\[
\kappa^{\kappa}\geq 2^{\kappa}>\kappa.
\]
Thus the displayed equality cannot hold in the intended infinite-cardinal setting, and the referenced theorem and its intended cardinal-arithmetic hypothesis must be identified before a meaningful necessity question can be stated. No replacement exponent is inferred here.

\item For singular \(\kappa\), that is, an infinite cardinal with \(\operatorname{cf}(\kappa)<\kappa\), determine the descriptive-set-theoretic complexity of embeddability on the intended space \(\operatorname{Mod}_{\kappa}(\varphi)\). Can this embeddability relation consistently be strongly invariantly universal for the intended \(\kappa\)-analytic quasi-orders? The sentence, coding, topology, embedding convention, meaning of strong invariant universality, and background theory for the consistency assertion remain unspecified.

\item In the corresponding universality theorems for combinatorial trees and generalized trees, can the tree classes be replaced by linear orders or groups while retaining the respective universality conclusions, including strong invariant universality where asserted? The particular theorems, cardinal assumptions, and structure-preserving maps for these replacements must be supplied to make this question precise.
\end{enumerate}

These are statement-only questions attributed to the supplied record; the missing specifications and incompatible printed hypothesis prevent a fully precise standalone formulation.
\subsection{Short English statement}
The record asks whether a cardinal-arithmetic assumption is necessary for strong invariant universality of tree embeddability, what happens at singular cardinals, and whether linear orders or groups can replace trees. The printed assumption \(\kappa^{\kappa}=\kappa\) is impossible for infinite \(\kappa\), and several essential definitions and theorem hypotheses are absent. The intended questions therefore require source clarification.
\subsection{Sources}
\begin{itemize}
\item Supplied local UnsolvedMath registry record, permanent ID \(30002479\), title \emph{Strong Invariant Universality at Singular Cardinals}, collection \emph{Oberwolfach Reports Open Problems}, exact source problem code \texttt{catalogue entry}. The record assigns the year 2014 and does not name a proposer.
\item Supplied report attribution: \emph{Set Theory} (2014), Oberwolfach Reports, pp. 91--144, DOI 10.4171/owr/2014/02. \url{https://doi.org/10.4171/owr/2014/02}. The supplied background describes Oberwolfach Reports as records of mathematical workshops and their open-problem discussions.
\item Supplied literature source: \emph{Oberwolfach set theory report}. \url{https://ems.press/content/serial-article-files/46488}.
\item Supplied literature source: Calderoni--Mildenberger--Motto Ros, ``universality'' (label supplied by the registry). \url{https://arxiv.org/abs/1903.08091}.
\item Literature assessment dated exactly \texttt{2026-08-21}: the supplied curation labels the problem partially solved and reports strong invariant universality and related reductions to generalized trees, graphs, and groups under cardinal-arithmetic hypotheses, while describing necessity, singular-cardinal complexity, and the full replacement questions as unsettled. This is supplied dated curation derived from the OWR material, not an independently verified literature assessment. The record explicitly requests independent verification of time-sensitive status and cited claims.
\end{itemize}
\end{document}
